\scanpage{177}
% archive foliation: 176
Pour prouver (ii), on est ramené à prouver que tout élément de $\mathrm{Prépic}'(X'/S')$ dont les images sont les mêmes, provient d'une classe de $\mathrm{Prépic}'(X/S)$. À cause de l'unicité, on voit que c'est local sur $S$, on peut donc supposer qu'on a une section de $X$ sur $S$. On peut alors identifier les \deleted{\unsure{classes}}\inserted{élts} de $\mathrm{Prépic}(X/S)$ à des faisceaux inversibles \inserted{(à isom. unique près)} sur $X$ à section marquée sur $S$, ce qui permet d'appliquer la théorie de la descente \uline{fid. plate}.

\uline{Corollaire 2} Même conclusion que Prop 1. (ii) pour descendre des \uline{faisceaux} sur $X$.

\uline{Corollaire 3} Supposons que la relation $\underline H^0_{X/S}(\underline O_X)=\underline O_S$ reste vraie par tte extension de la base (i.e., $X$ coh. plat sur $S$ en dim 0, et $\underline H^0_{X/S}(\underline O_X)=\underline O_S$), \inserted{et $X/S$ a localement une section} Alors le foncteur $\underline{\underline{\mathrm{Prépic}}}'_{X/S}$ est compatible avec [= ex. et pour] la descente fidèlement plate quasi-compacte.

